\documentclass[11pt]{article}
\usepackage[margin=0.78in]{geometry}
\usepackage{amsmath,amssymb,booktabs,enumitem,microtype}
\usepackage[hidelinks]{hyperref}
\setlist{nosep,leftmargin=*}
\setlength{\parindent}{0pt}
\setlength{\parskip}{0.45em}
\newcommand{\code}[1]{\texttt{#1}}
\newcommand{\file}[1]{\texttt{game\_PIRS/#1}}

\title{How the Policy Interest Rate Simulator Runs\\\large A Mechanistic Code Walkthrough}
\author{}
\date{}

\begin{document}
\maketitle
\vspace{-2.5em}

\section{The executable path and the two nested loops}
The current web entry point is \file{app.py}.  Running
\code{streamlit run game\_PIRS/app.py} executes the module from top to bottom,
defines its helper functions, and calls \code{main()} under the usual
\code{\_\_main\_\_} guard.  Streamlit then reruns that script after widget
interactions.  Persistent objects therefore live in \code{st.session\_state},
not in ordinary local variables.  The older \file{executable.py} is a separate
Tkinter front end; it drives the same \code{Economy} object but is not part of
the Streamlit path described below.

There are two important loops.  The first is a hidden, forty-quarter bootstrap
inside \code{\_new\_game}.  The second is the visible policy loop, advanced one
quarter whenever the player submits an interest rate.

\subsection*{Loop 1: manufacture the inherited economy}
The start-page selections call \code{\_new\_game(difficulty, scenario,
mandate)}.  Mechanically, it does the following.
\begin{enumerate}
  \item It constructs \code{Economy(difficulty="central\_banker")}.  The
  constructor creates parameters, random indicators, a random initial nominal
  rate, reputation $0.8$, a random central-bank persona, an eight-slot event
  queue, history storage, and the plotting compatibility store.
  \item Scenario setup may overwrite a few indicators, select a bootstrap
  persona, suppress late events, or force a named event.  Notably,
  \code{\_sample\_scenario} currently returns \code{None}; most scenario
  behavior is imposed by the surrounding bootstrap code rather than passed to
  \code{Economy}.
  \item For forty iterations, the persona chooses a Taylor-style rate and then
  \code{simulate\_quarter()} advances the model.  For example, the ``good''
  persona requests
  \begin{equation}
    i_t^*=r_t^*+\pi_t+0.5(\pi_t-2)
       +0.5(u_t^n-u_t),
  \end{equation}
  after which the rate is rounded to a quarter point and bounded below by zero.
  Different personas change the intercept, targets, and inflation/employment
  weights.
  \item Only after the bootstrap does \code{set\_difficulty} apply the player's
  selected shock scale and event cooldown.  Thus inherited history is always
  generated under central-banker settings; the visible term uses the selected
  difficulty.
  \item The last ten unemployment observations determine the dual-mandate
  unemployment target (rounded to an integer).  The economy, counters, news,
  targets, and initial term readings are stored in session state.
\end{enumerate}

\subsection*{Loop 2: accept one policy decision}
On submission, \code{\_next\_quarter} writes the player's nominal rate directly
to \code{econ.interest\_rate}, calls \code{simulate\_quarter}, adds any selected
event to the news log, increments the visible term counters, and checks whether
sixteen policy quarters have elapsed.  Streamlit reruns and redraws charts from
the accumulated \code{Variables} histories.  In compact form,
\[
  \boxed{i_t\text{ chosen}}\longrightarrow
  \boxed{\text{event queue}}\longrightarrow
  \boxed{\varepsilon_t\text{ drawn}}\longrightarrow
  \boxed{\text{AD--AS solved}}\longrightarrow
  \boxed{u_t,\rho_t,\text{history}}.
\]

\newpage
\section{Which script owns which part of the run?}
The simulator is not one long script.  The UI calls a coordinator, and that
coordinator calls small modules in a fixed order.  The following map names the
specific responsibility of each file and the object or value it hands onward.

\begin{center}
\begin{tabular}{@{}p{0.23\linewidth}p{0.35\linewidth}p{0.34\linewidth}@{}}
\toprule
File & What it executes or defines & What leaves the file \\
\midrule
\file{app.py} & Builds the Streamlit screens, creates a game, runs the bootstrap,
accepts a rate, and requests one quarter. & A player rate enters
\code{Economy}; returned event text enters the news log. \\
\file{economy.py} & Defines \code{Economy}, the turn coordinator.  It calls the
event, shock, equation, reputation, and history layers in order. & A result
dictionary goes back to the UI; new state remains on the object. \\
\file{indicators.py} & Defines the mutable \code{EconomicIndicators} dataclass
and draws a random initial macroeconomic state. & Current inflation,
unemployment, natural unemployment, $r^*$, and growth. \\
\file{parameters.py} & Defines \code{EconomyParameters}; it contains all curve
coefficients, bounds, shock sizes and correlations, and solver controls. & One
parameter object read by shocks, drift, and laws of motion. \\
\file{events.py} & Defines \code{ProbTerm}, \code{GameEvent}, probability
formulas, descriptions, and eight-quarter effect schedules. & The event catalog
loaded once by \code{EventEngine}. \\
\file{event\_engine.py} & Chooses at most one new event, adds its schedule to the
queue, and consumes effects due now. & An \code{EventOutcome} containing name,
description, and a current-effect dictionary. \\
\file{shocks.py} & Converts standard deviations and correlations into a
covariance matrix and makes one multivariate-normal draw. & A four-number NumPy
array $[\varepsilon^\pi,\varepsilon^d,\varepsilon^{u^n},\varepsilon^{r^*}]$. \\
\file{laws\_of\_motion.py} & Defines AD, AS, Okun's law, the Newton solver, and
\code{solve\_ad\_as}. & One immutable \code{MotionResult} with six solved or
derived numbers. \\
\file{reputation.py} & Applies transparent bonuses and penalties after the new
inflation and unemployment are known. & A clipped reputation number in $[0,1]$. \\
\file{history.py} & Defines immutable quarterly rows and builds the event
engine's history-shaped input. & An appended \code{HistoryEntry} and, on demand,
a dictionary of historical series. \\
\file{variables.py} & Maintains the smaller legacy current-value/history view
used by both interfaces. & Lists read by charts and labels. \\
\file{endgame\_logic.py} & Scores the last term and classifies the central
banker's public style. & The text displayed in the end-of-term dialog. \\
\bottomrule
\end{tabular}
\end{center}

Two other files are presentation alternatives rather than extra stages in a
Streamlit turn.  \file{executable.py} and the larger \file{main\_gui3.py} are
Tkinter interfaces that also call \code{Economy}; \file{visualization.py} is a
small Matplotlib helper.  \file{utils.py} supplies the model's linear real-rate
formula.
Its older \code{generate\_shocks} helper is not the one imported by
\file{economy.py}; the live coordinator imports from \file{shocks.py}.

\subsection*{The concrete call stack for a click}
When the player submits a rate, the main path is
\begin{align*}
 \file{app.py}:\ &\code{\_next\_quarter(rate)}\\
 &\downarrow\ \code{Economy.adjust\_interest\_rate(rate)}\\
 &\downarrow\ \code{Economy.simulate\_quarter()}\\
 \file{economy.py}:\ &\code{EventEngine.advance()}
 \rightarrow\code{generate\_shocks()}
 \rightarrow\code{solve\_ad\_as()}\\
 &\rightarrow\code{update\_reputation()}
 \rightarrow\code{EconomicHistory.append()}
 \rightarrow\code{Variables.update()}.
\end{align*}
The UI never calls AD, AS, or Okun's law directly.  Conversely, the equation
module knows nothing about Streamlit, events, history, or scoring.  The
\code{Economy} coordinator is the bridge between those layers.

\newpage
\section{One call to \code{simulate\_quarter}, instruction by instruction}
\file{economy.py} is the coordinator.  It deliberately keeps event selection,
random shocks, equations, reputation, and storage in separate modules.  One
call advances exactly one model quarter in this order.

\subsection*{1. Build the information set}
\code{EconomicHistory.event\_snapshot} converts the append-only typed history
into the legacy dictionary expected by event probability functions: complete
series for inflation, rates, unemployment, reputation, the recent event lists,
and the user-facing quarter.  Event probabilities therefore see state only up
to the start of this quarter.

\subsection*{2. Select, enqueue, and consume an event}
\file{event\_engine.py} first checks a special cooperative fiscal rule in
principles mode, then the difficulty cooldown.  Otherwise every eligible event
$j$ computes its history-dependent probability $p_{j,t}$ and receives an
independent uniform draw:
\begin{equation}
  j\text{ is a candidate}\quad\Longleftrightarrow\quad
  U_{j,t}<\min(1,\max(0,p_{j,t})),\qquad U_{j,t}\sim U(0,1).
\end{equation}
If several events pass, one candidate is chosen uniformly, so at most one new
event is announced.  Its schedules are added element-by-element to the shared
eight-quarter queue.  The engine immediately pops queue slot zero and appends
an empty tail slot.  Consequently, a new event's first scheduled effect occurs
now, while its later effects coexist additively with other active schedules.

Suppose an event schedules inflation effects $(1,0.5,0,\ldots)$ and an existing
event already placed $0.2$ in next quarter's slot.  This call consumes $1$; the
next call consumes $0.5+0.2$.  The queue is thus a small convolution-like
accumulator, not a list of event objects.

The coordinator treats inflation specially.  It removes the key
\code{inflation} from current effects and saves it as $e_t^{\pi}$.  Other event
effects immediately shift the state:
\[
 i_t\mathrel{+}=e_t^i,\quad r_t^*\mathrel{+}=e_t^{r^*},\quad
 u_{t-1}\mathrel{+}=e_t^u,\quad u_t^n\mathrel{+}=e_t^{u^n}.
\]
The inflation event is instead fed into aggregate supply during the simultaneous
AD--AS solve.  This distinction matters: direct unemployment effects alter the
base from which Okun's law runs; inflation effects alter the curve itself.

\subsection*{3. Generate background shocks}
\file{shocks.py} draws four correlated Gaussian innovations in the order
inflation, demand, natural unemployment, and equilibrium real rate:
\begin{align}
 \boldsymbol\varepsilon_t&\sim\mathcal N(\mathbf0,\Sigma),&
 \Sigma&=D R D,&D&=\operatorname{diag}(s\boldsymbol\sigma).
\end{align}
$R$ and $\boldsymbol\sigma=(0.3,0.2,0.05,0.1)$ come from
\file{parameters.py}; $s$ is $0.1$, $0.5$, or $1$ in principles, senior, or
central-banker mode.  Before AD--AS, natural unemployment and $r^*$ both
mean-revert slowly toward their long-run anchors:
\begin{align}
 u_t^n&=\max\{2,\,u_{t-1}^n-0.02(u_{t-1}^n-5)+\varepsilon_t^{u^n}\},\\
 r_t^*&=r_{t-1}^*-0.02(r_{t-1}^*-0.5)+\varepsilon_t^{r^*}.
\end{align}

\newpage
\section{The equation engine: from the rate to inflation and jobs}
\file{laws\_of\_motion.py} solves two equations jointly.  Let $g_t$ be GDP
growth, $g^p=2$ potential growth, $i_t$ the chosen nominal rate, and $r_t^*$ the
equilibrium real rate.  Aggregate supply (a Phillips curve) is
\begin{equation}
 \pi_t=\underbrace{2}_{\pi^e}
       +\underbrace{0.25}_{\beta_\pi}(g_t-g^p)
       +\varepsilon_t^\pi+e_t^\pi. \tag{AS}
\end{equation}
Aggregate demand is
\begin{equation}
 g_t=\underbrace{2}_{\alpha}
     +\underbrace{(-0.5)}_{\beta_y}
       (i_t-\pi_t-r_t^*)+\varepsilon_t^d. \tag{AD}
\end{equation}
This AD equation uses $i_t-\pi_t$ as a linear approximation to the real rate.
Because inflation appears in AD and output appears in AS, neither equation is
run ``first.''  The code searches for their intersection.

For candidate vector $x=(\pi,g)^\top$, it constructs the residual
\begin{equation}
 F(x)=\begin{bmatrix}
 \pi-AS(g)\\ g-AD(\pi)
 \end{bmatrix}.
\end{equation}
Starting from $(2,2)$, Newton iterations continue until
$\|F(x)\|_\infty\le10^{-9}$ or fifty iterations have run:
\begin{align}
 J_{:,k}(x)&\approx\frac{F(x+h e_k)-F(x-h e_k)}{2h},
 &h&=10^{-5},\\
 J(x)\Delta x&=-F(x),&x&\leftarrow x+\Delta x.
\end{align}
Although today's equations are linear, the numerical Jacobian allows a future
nonlinear curve to use the same solver.  A singular Jacobian or failure to
converge raises \code{NumericalSolutionError}; the coordinator does not swallow
that error.

Only after the intersection does Okun's law update unemployment:
\begin{equation}
 u_t=\operatorname{clip}_{[1,99]}
       \left(u_{t-1}-0.7(g_t-g^p)\right). \tag{Okun}
\end{equation}
Notice the precise wiring: current natural unemployment is recorded and is used
by automated personas, but it is not an argument to this Okun equation.  A
positive output gap lowers unemployment relative to its previous value; a
negative gap raises it.

\subsection*{Commit and feedback}
The returned \code{MotionResult} overwrites inflation, GDP growth, and
unemployment.  For display and reputation, the code uses the same linear real
rate approximation used by aggregate demand:
\begin{equation}
 r_t^{\mathrm{real}}=i_t-\pi_t.
\end{equation}
Reputation then receives additive rule-based changes (for example, falling
inflation adds $0.02$, inflation above six subtracts $0.05$) and is clipped to
$[0,1]$.  This reputation is recorded and can affect later event probabilities,
closing a feedback path from outcomes to future event risk.

\newpage
\section{Exactly when \code{MotionResult} appears}
\code{MotionResult} is a frozen dataclass declared near the top of
\file{laws\_of\_motion.py}.  It is neither the whole economy nor a history row.
It is a short-lived return package for the answer to this quarter's equation
problem.  Its six fields are
\begin{center}
\begin{tabular}{@{}p{0.25\linewidth}p{0.66\linewidth}@{}}
\toprule
Field & Meaning and source \\
\midrule
\code{inflation} & The $\pi_t$ coordinate of the numerical AD--AS intersection. \\
\code{output\_growth} & The $g_t$ coordinate of that same intersection. \\
\code{unemployment} & The bounded $u_t$ produced when Okun's law uses
\code{output\_growth} and last quarter's unemployment. \\
\code{output\_gap} & The derived difference $g_t-g^p$. \\
\code{aggregate\_demand} & AD evaluated again at the solved inflation.  At a
valid intersection it equals \code{output\_growth}, up to tolerance. \\
\code{aggregate\_supply} & AS evaluated again at the solved growth.  At a valid
intersection it equals \code{inflation}, up to tolerance. \\
\bottomrule
\end{tabular}
\end{center}

It comes into existence only at the final \code{return MotionResult(...)} in
\code{solve\_ad\_as}.  Before that line, the Newton solver works with an ordinary
two-element NumPy candidate array.  After that line, control returns to
\code{Economy.simulate\_quarter}, where the local variable is literally named
\code{motion}:
\begin{equation*}
 \underbrace{\code{motion = solve\_ad\_as(...)}}_{\text{create and receive one MotionResult}}
 \quad\longrightarrow\quad
 \begin{cases}
  \code{\_commit\_motion(motion, previous\_inflation)},\\
  \code{\_record\_quarter(motion, shocks, event\_name)}.
 \end{cases}
\end{equation*}

\subsection*{First consumer: commit the new live state}
\code{\_commit\_motion} reads three fields.  It copies
\code{motion.inflation} into \code{indicators.inflation\_rate},
\code{motion.output\_growth} into \code{indicators.gdp\_growth}, and
\code{motion.unemployment} into \code{indicators.unemployment\_rate}.  The
inflation copy is bounded below by
\code{parameters.minimum\_inflation}; the bound is a configurable model
parameter rather than a literal in the coordinator.  The object has now changed
from the start-of-quarter state to the end-of-quarter state.  The method then
computes $i_t-\pi_t$ from the newly committed inflation and updates reputation.
It does not keep the \code{MotionResult} as an attribute.

\subsection*{Second consumer: write an auditable row}
\code{\_record\_quarter} reads the other three fields as well.  It stores
\code{motion.aggregate\_demand}, \code{motion.aggregate\_supply}, and the output
gap alongside the committed indicators, shocks, event, rates, and reputation.
The output gap is also returned to the UI under the legacy key
\code{gap\_effect}.  Then \code{\_update\_variables} copies selected live values
to chart histories.  Once these calls return, no reference to that
\code{MotionResult} remains; the durable versions of its values are the current
indicators and the new history row.

This small object prevents a half-solved economy from leaking out of the
equation module.  AD, AS, and Okun calculations finish first; only then does the
coordinator commit all three headline outcomes.  Its frozen status also means a
consumer cannot accidentally rewrite the solver's answer.  The component tests
use the redundant curve fields to assert
\[
 \code{output\_growth}=\code{aggregate\_demand},\qquad
 \code{inflation}=\code{aggregate\_supply},
\]
which checks that the returned point really lies on both curves.

\newpage
\section{Storage, presentation, and termination}
Every completed quarter becomes an immutable \code{HistoryEntry}.  It stores
the solved outcomes, both curve values, all four innovations, the event name,
policy and real rates, natural unemployment, $r^*$, and reputation.  This is the
authoritative, rich record used to construct future event information.

The same state is also copied into \code{Variables}, a narrower compatibility
store used by existing charts.  Each \code{update(name,value)} both replaces the
latest value and appends to that variable's list.  Therefore a chart redraw is
read-only: it obtains the inflation, unemployment, interest-rate, and sometimes
natural-unemployment histories, reshapes them for Altair, and overlays the
player-start and optional mandate-target lines.  It does not advance the model.

The coordinator finally increments \code{current\_quarter}.  This explains a
common off-by-one trap: the history entry is labeled with the quarter just
solved, then the public counter points at the next quarter.  The result returned
to the UI contains the new event's description and name, the solved output gap,
and the four background shocks.  The returned shock vector does not include the
event inflation addition, although the history's \code{inflation\_shock} does.

After the sixteenth player decision, \code{\_finish\_game\_if\_needed} slices
the latest sixteen history observations and calls \file{endgame\_logic.py}.
Scoring uses the last twelve term observations.  Inflation targets two percent;
the dual mandate additionally uses the inherited unemployment target.  Average
inflation within one point (and, for the dual mandate, unemployment no more
than one point above target) is provisionally successful.  Three or more
extreme misses override success.  Otherwise recent improvement can produce a
mixed verdict.  The public persona classification counts patterns in inflation,
unemployment, and real rates, and the final dialog names economically active
term events.

\subsection*{Dependency and feedback map}
\begin{center}
\begin{tabular}{@{}p{0.20\linewidth}p{0.31\linewidth}p{0.40\linewidth}@{}}
\toprule
Producer & Immediate consumer & Later feedback \\
\midrule
Player/persona rate & AD real-rate gap & outcomes, charts, event history \\
Event probability & eight-slot effect queue & event cooldown and recent-event rules \\
Shock generator & drift equations and AD--AS & stored shock history \\
AD--AS intersection & Okun, commit step & next quarter's state \\
Committed outcome & reputation and history & future event probabilities \\
History/Variables & events, scoring, charts & player observes and chooses next rate \\
\bottomrule
\end{tabular}
\end{center}

For reproducible experiments, seed NumPy before constructing the economy;
selection, personas, initial conditions, and shocks all use NumPy's random
machinery.  Tests can also inject an \code{initial\_state} and custom
\code{EconomyParameters}.  The component tests verify the AD--AS intersection,
the direction of rate effects, history growth, and one-slot-at-a-time event
schedule consumption.

\newpage
\section{Summary (200 words)}
The game starts in \file{app.py}. It creates an economy and runs forty hidden
quarters. Then a computer central banker chooses the interest
rate. The player then takes control. Each time the player enters a rate,
\file{economy.py} runs one quarter.

First, \file{event\_engine.py} may choose an event. It puts the event's effects
into a queue. Effects happen now or in later quarters.
Next, \file{shocks.py} makes four random shocks. These shocks can change
inflation, demand, normal unemployment, and the neutral real interest rate.

Then \file{laws\_of\_motion.py} solves the equations. The interest rate
affects demand. Demand and supply together set inflation and growth. Okun's law
uses growth to set unemployment. The script puts these six answers into an
object called \code{MotionResult}. \file{economy.py} reads it, updates
the live economy, and saves the results. \file{history.py} keeps the full record.
\file{variables.py} keeps the lists used by the charts.

The new results affect what can happen next. They can change future event
chances, reputation, and the player's next choice. After sixteen player
quarters, \file{endgame\_logic.py} checks inflation, unemployment, interest
rates, and events. It then tells the player whether the policy goal was met and
how the public remembers the central banker.

\end{document}
